\section{Experimental Setup}
\label{sec:experimental_setup}

We organize the evaluation around three questions, each addressed by a
predeclared sub-study: \textbf{(Q1)} does whole-pool prompting move the
Pareto frontier on the comparisons-axis and wall-clock-axis at
matched-\texttt{hits} cost? \textbf{(Q2)} which mechanism --- joint
elicitation, whole-pool selection, or single-extreme selection --- is
responsible for any movement? \textbf{(Q3)} how does whole-pool prompting
behave under input-order perturbations and varying passage lengths? The
matrix below is staged so that order-stability and per-query parse
reliability are validated before the bulk of compute is committed.

\subsection{Datasets and First-Stage Retrieval}

We evaluate on the standard zero-shot LLM reranking benchmarks: \textbf{TREC
Deep Learning 2019}~\citep{craswell2020dl19} and \textbf{TREC Deep Learning
2020}~\citep{craswell2021dl20}, both built on the MS MARCO passage
collection. DL19 contains 43 judged queries with graded relevance, DL20
contains 54. First-stage retrieval is BM25~\citep{robertson2009bm25} as
implemented in Pyserini~\citep{lin2021pyserini}, retrieving the top-$N$
candidates per query for $N \in \{10,20,30,40,50\}$. Matched-\texttt{hits}
discipline (see \S\ref{ssec:baselines}) requires that every reranking method
is evaluated against a baseline that re-ranked the \emph{same} BM25 top-$N$
list, since archived runs at $\texttt{hits}{=}100$ cannot be subsetted to
$\texttt{hits}{=}50$ without changing the candidate pool that the reranker
ever sees.\footnote{The reference run-script reads only the top-\texttt{hits}
BM25 candidates before any LLM call, so a $\texttt{hits}{=}50$ rerun is a
genuinely different experiment from a $\texttt{hits}{=}100$ run with
post-hoc top-50 truncation.} All experiments report nDCG@10 (primary). We
also log nDCG@100, MAP@10, MAP@100, and Recall@1000 against the official
qrels via \texttt{trec\_eval}; secondary metrics are reported in the
appendix unless space permits inclusion in the main tables.

\subsection{Models}

We instantiate MaxContext on six instruction-tuned Qwen
checkpoints~\citep{qwen3technical}: \textbf{Qwen3-4B}, \textbf{Qwen3-8B},
\textbf{Qwen3-14B}, \textbf{Qwen3.5-4B}, \textbf{Qwen3.5-9B}, and
\textbf{Qwen3.5-27B}. We restrict to Qwen because only the
generation-scoring code path on Qwen3 / Qwen3.5 in our implementation
exercises true joint best-worst elicitation as a single decoding step;
encoder-decoder backbones such as Flan-T5~\citep{chung2024flant5} and the
likelihood-scoring path collapse joint elicitation to a best-only proxy via
argmax / argmin over independent label likelihoods. \textbf{All MaxContext
claims in this paper are therefore scoped to Qwen-generation models; we do
not claim model-family generality.} Llama- and Mistral-based extension is
conditional on family-specific sanity-smoke validation (parse stability at
$N{=}50$, order robustness, refusal-pattern coverage) and is out of scope
for this submission.

Two contextual notes apply across the model lineup. \textbf{Native context
length:} Qwen3-{4B,8B,14B} ship with $32{,}768$-token native context;
Qwen3.5-9B and the larger Qwen3.5 checkpoints ship with $262{,}144$-token
native context. At passage length $512$, a $50$-document pool plus query,
prompt scaffolding, and decoding budget consumes well under $30$K tokens,
so no model in our lineup is ever close to its native limit. \textbf{Thinking
mode disabled:} Qwen3 / Qwen3.5 generation runs use
\texttt{enable\_thinking=False}, capping the decoding budget at $\sim$128
output tokens to discourage long chain-of-thought and keep the parse target
to a structured \texttt{Best: <int>, Worst: <int>} (or single-int) string.

\subsection{MaxContext Configuration}
\label{ssec:max_config}

All MaxContext runs use a fixed harness with hard invariants. Each ranker
asserts at construction time that
\texttt{pool\_size} $=$ \texttt{hits} $=$ \texttt{ranker.k}, that
\texttt{num\_permutation} $=$ $1$, that the model family is in the Qwen3 /
Qwen3.5 allowlist, and that the runtime-rendered prompt fits inside the
model's context budget with $128$ tokens reserved for decoding. Truncation
is forbidden: any prompt that would exceed the budget aborts the run with a
diagnostic. Document labels are numeric ($1$..$N$); the parser tightens the
upstream regex to validate that extracted integers lie in $[1, N]$ and to
reject duplicate or out-of-range labels with no silent default. Out-of-range
labels and duplicate best/worst labels trigger an immediate run-abort
rather than a fallback ranking. For the single-extreme variants, an
$n_{\text{docs}}{=}2$ deterministic BM25 endgame resolves the last pair:
\textsc{MaxContext-TopDown} keeps the higher-BM25 document at the next top
rank (so the BM25 endgame decides ranks $N{-}1$ vs.\ $N$, the tail of the
ranking), while \textsc{MaxContext-BottomUp} keeps the lower-BM25 document
at the next bottom rank (so the BM25 endgame decides ranks $1$ vs.\ $2$,
the head of the ranking). Because nDCG@10 is sensitive to head-of-ranking
decisions, we report \textsc{TopDown} and \textsc{BottomUp} on the Pareto
plot \emph{separately} and decline any causal claim about whether the
single-extreme bypass position causes the observed quality gap.

\subsection{Baselines}
\label{ssec:baselines}

\paragraph{Matched-\texttt{hits} setwise / DualEnd lineage.}
For every \texttt{hits}\,$=$\,pool-size value $N \in \{10, 30, 50\}$
(the cheapest, mid-range, and deepest anchors of our pool sweep), we
re-rerank the same BM25 top-$N$ list with four predeclared baselines:
\textbf{TopDown-Heap (TD-Heap)} and \textbf{TopDown-Bubble (TD-Bubble)}
(both upstream setwise sorts at \texttt{num\_child}$=$$3$);
\textbf{DualEnd-Cocktail (DE-Cocktail)} (cocktail-shaker sort over a joint
best-worst comparator); and \textbf{DualEnd-Selection (DE-Selection)}
(double-ended selection sort over the same comparator). All four use the
public setwise reference codebase~\citep{ielab_llmrankers}, our DualEnd
extension built on top, and the same six Qwen models, two datasets, and
passage length $512$ as the MaxContext primary runs. If Phase 2's
order-stability gate (see below) fails at $N{=}50$ and forces the production
matrix to a smaller pool size, the matched-\texttt{hits} baseline grid is
extended to include the chosen smaller $N$ ($N{=}20$ by default) so that the
non-inferiority rule remains evaluated against same-\texttt{hits} baselines.

The headline rule for \textbf{(Q1)} is non-inferiority: at matched
\texttt{hits} $N$, \textsc{MaxContext-DualEnd} must be no worse than the
\emph{best of} \{DE-Cocktail, DE-Selection, TD-Bubble\} on nDCG@10 with a
predeclared margin of $\Delta \geq -0.003$ (point estimate) and a
bootstrap-CI lower bound of $\geq -0.005$, while reducing the average
per-query LLM call count. We use the best-of rule rather than a single
baseline because DE-Selection holds the only Bonferroni-significant
DualEnd win in our completed small-window experiments; comparing only
against DE-Cocktail would effectively cherry-pick an easier baseline. If a
single-extreme variant rather than DualEnd dominates the Phase 4 frontier,
the same non-inferiority rule is reported for that variant against
\emph{same-direction} matched-\texttt{hits} baselines (TD-Heap or TD-Bubble
for \textsc{MaxContext-TopDown}; the BU-Heap upstream baseline for
\textsc{MaxContext-BottomUp}), so the headline framing scales with whichever
mechanism is empirically responsible for the frontier movement.

\paragraph{Why TourRank is not a primary baseline.}
Our primary comparisons are within the matched-\texttt{hits} setwise /
DualEnd family because the main claim concerns whole-pool prompting under
controlled pool depth. TourRank~\citep{chen2025tourrank} varies different
budget knobs --- number of tournaments, group size, advancement fraction,
and parallel tournament schedule --- so a fair head-to-head requires a
separate matched-budget sweep. We report TourRank as related work
(\S\ref{ssec:rw_extensions}) and leave a full MaxContext--TourRank
comparison to future work. \textbf{TourRank is therefore not part of the
non-inferiority rule above; the matched-\texttt{hits} non-inferiority claim
is internal to the setwise / DualEnd family.}

\paragraph{Top-$m$ + DualEnd cascade.}
The complementary operating point is evaluated under the same Qwen lineup,
datasets, and passage length. The cascade starts from the BM25 top-$100$
candidates --- a deeper pool than the matched-\texttt{hits} setwise
baselines, since the cascade's role is to probe whether prefiltering with
Top-$m$ selection lets us refine over a larger initial pool than MaxContext
alone can fit in context. Hyperparameters are predeclared as $G \in \{10,
20\}$ (setwise group size at the Top-$m$ stage), $m{=}5$, $r \in \{1, 3\}$
(repeated stratified groupings), and a finalist pool capped at $|F|{=}30$
that is refined directly with \textsc{MaxContext-DualEnd}
($\lfloor |F|/2 \rfloor$ dependent calls; no sharding required at
$|F|{\le}50$). The finalist scoring rule is predeclared:
\[
\mathrm{score}(d) = \mathrm{vote}^+(d) - \lambda \,\mathrm{vote}^-(d) + \beta \,\mathrm{prior}_{\mathrm{BM25}}(d),
\]
with $\lambda{=}0.25$, $\beta{=}0.10$, and deterministic tie-breaks (higher
positive vote rate first, then better BM25 rank, then stable document ID).
\textit{Per-query normalization:} $\mathrm{vote}^+(d)$,
$\mathrm{vote}^-(d)$, and $\mathrm{prior}_{\mathrm{BM25}}(d)$ are each
min--max scaled to $[0,1]$ within the query before being combined, so the
$\lambda$ / $\beta$ hyperparameters operate on a comparable scale across
queries. Cascade hyperparameters are reported as a fixed sensitivity check
rather than as a tuned competing method, and finalist-pool recall (the
fraction of qrels-relevant documents retained in $F$) is reported
separately as a diagnostic of how aggressively the Top-$m$ stage prunes.

\subsection{Cost Reporting and Statistics}

\paragraph{Five cost axes.}
Following~\citet{peng2025flopsreranking}, every reported method has all
five axes recorded and figured: (i) average per-query \emph{dependent} LLM
calls, labelled ``Avg comparisons'' for parity with prior setwise tables;
(ii) average per-query prompt tokens; (iii) average per-query completion
tokens; (iv) average per-query total tokens; (v) end-to-end wall-clock
latency, measured on identical NVIDIA H100 GPUs via Vast.ai. We also log
per-call parse status (\texttt{full\_parse}, \texttt{fallback\_parse}, or
\texttt{abort}) and report the \emph{parse-fallback rate} as the fraction
of LLM calls per query with a \texttt{fallback\_parse}, averaged over the
test set. \texttt{fallback\_parse} indicates the parser used a secondary
extractor (e.g., a numeric-only refusal regex) and recovered a valid label
within $[1, N]$; an \texttt{abort} indicates an out-of-range label,
duplicate best/worst labels, or refusal that the parser cannot recover from
--- aborts halt the entire run rather than substitute a fallback ranking. \textbf{We do not claim token
efficiency:} on prompt-token and total-token axes, MaxContext is expected
to be more expensive than small-window setwise; the only legitimate
efficiency claims are on the comparisons-axis and wall-clock-axis.

\paragraph{Significance testing.}
Pairwise differences in per-query nDCG@10 are tested with paired
approximate-randomization tests. Bonferroni correction is applied across the
within-family comparison set (number of model-dataset configurations per
family, typically 12 within Qwen). Bootstrap 95\% confidence intervals on
the per-query nDCG@10 difference are reported alongside point-estimate
deltas, so non-inferiority claims combine point and interval evidence per
\S\ref{ssec:baselines}.

\subsection{Staged Matrix and Launch Gates}

The full MaxContext production matrix runs on a shared SLURM cluster. Total
sizing is approximately \textbf{420 runs} ($\sim$$840$ H100 hours of
compute) across DualEnd, TopDown, and BottomUp pool sweeps plus the
matched-\texttt{hits} baseline grid; the mandatory pre-production gate
(Phases 2 and 3 below) adds $\sim$$56$ H100 hours, and the contingent
smaller-pool fallback adds another $\sim$$36$ H100 hours. Detailed launchers
and command lines are released with the code.

\paragraph{Phase 1 -- Unit sanity.}
Single-GPU smoke test on Qwen3-4B, DL19, $N \in \{10,20,30,40,50\}$, $pl=512$.
Pass criteria: zero parse aborts, zero truncations, output is a permutation
of the input pool. Phase 1 has already passed for all three MaxContext
variants in our development runs.

\paragraph{Phase 2 -- Order-robustness pilot (Q3 launch gate).}
On Qwen3-4B and Qwen3.5-9B, DL19 and DL20, at $N{=}50$ and $pl{=}512$, we
run \textsc{MaxContext-DualEnd} with three input orderings: BM25-forward (as
retrieved), BM25-inverse, and a fixed-seed random shuffle (seed $\!=\!42$
for reproducibility). \textbf{Pass criterion (primary):} maximum pairwise
nDCG@10 difference across the three orderings must be $\leq 0.01$ for both
models on both datasets. \textbf{Phase 2 also produces three secondary
position-bias-at-scale diagnostics} that are reported regardless of pass /
fail:
\begin{itemize}
\item \textit{Selected-position concentration:} per-call frequency at which
each in-prompt position $i \in [1, N]$ is selected as best (and as worst,
under DualEnd). Concentration $> 30\%$ at any single position under
forward ordering is flagged.
\item \textit{Selected-doc stability:} fraction of queries on which the
same docid is chosen between forward and inverse orderings; $< 50\%$ is
flagged as order-dependent.
\item \textit{Parse-fallback location distribution:} the in-prompt
position(s) most often associated with parse-fallback events, used to
inform parser hardening before Phase 4.
\end{itemize}
Failing the primary gate at $N{=}50$ re-targets the production matrix to
the contingent fallback at $N{=}20$ (where context-attention degradation is
less likely); persistent failure at $N{=}20$ aborts the pivot. Position-
bias-at-scale results are reported separately from the $w{=}4$ letter-
scheme analyses in our prior work, since small-window evidence does not
automatically extend to $N{=}50$.

\paragraph{Phase 3 -- Matched-\texttt{hits} non-inferiority (Q1 launch gate).}
Four \texttt{(model, dataset)} pairs (Qwen3-8B and Qwen3.5-9B $\times$ DL19
and DL20) compare \textsc{MaxContext-DualEnd} at $N{=}50$ against
\{DE-Cocktail, DE-Selection, TD-Bubble\} at matched \texttt{hits}{=}$50$
under the best-of non-inferiority rule of \S\ref{ssec:baselines}. The
DE-Cocktail and DE-Selection runs are launched fresh for the four pairs;
the TD-Bubble runs at $\texttt{hits}{=}50$ are part of the same Phase 4B
matched-\texttt{hits} baseline grid (and so are not double-counted in the
Phase 3 gate's compute budget unless the corresponding TD-Bubble run is
not yet on disk). \textbf{Pass criterion:} non-inferiority on at least 3
of 4 pairs. Phase 3 is a \emph{launch gate}, not a final generalization
claim: a 4-pair test on the small TREC DL test sets is large enough to
catch model-specific or dataset-specific quirks that would otherwise
contaminate Phase 4, but the final evidence must come from the full Phase
4 6-model $\times$ 2-dataset production matrix with bootstrap CIs and
paired tests across all $12$ configurations.

\paragraph{Phase 4 -- Production sweeps (Q1 + Q2).}
Conditional on Phases 2--3 passing. Three sweeps:
(i) Study A --- pool-size sweep $N \in \{10,20,30,40,50\}$ for
\textsc{MaxContext-DualEnd} on six Qwen models $\times$ two datasets ($60$
runs);
(ii) Phase 4D --- the same sweep for \textsc{MaxContext-TopDown} ($60$ runs);
(iii) Phase 4E --- for \textsc{MaxContext-BottomUp} ($60$ runs);
plus the Phase 4B matched-\texttt{hits} baseline grid at $N \in \{10,30,50\}$
(or $\{10,20,30,50\}$ if Phase 2 forced the smaller-pool fallback) across
four sort families ($144$ runs at the default grid). Phase 4 produces the
headline Pareto plots and the three-variant mechanism ablation.

\paragraph{Phase 5 -- Passage-length sensitivity (Q3).}
At the predeclared production pool size from Phase 4, the passage-length
sweep $pl \in \{64, 128, 256, 512\}$ runs both
\textsc{MaxContext-DualEnd} (treatment, $48$ runs) and a DE-Selection-with-
\texttt{num\_child}$=$$3$ control arm at the same $pl$ ($48$ runs).
Confound rule: if shorter $pl$ wins under MaxContext but not under the
control, attribute to long-context attention degradation; if shorter $pl$
wins under both, attribute to passage-length sensitivity in the underlying
Qwen scoring; any other pattern triggers further investigation before any
attribution.

\subsection{Reproducibility}

All MaxContext code, prompt templates, parser, and launcher scripts are
released under the same open-source licence as the underlying setwise
codebase. Each completed run emits a JSON Lines log containing per-query
LLM-call traces (prompt, response, parse status, label scheme), the
deterministic seeds used for input-ordering pilots, the SLURM job ID, and
the model checkpoint hash, so that downstream analyses --- in particular
the position-bias-at-scale stratification --- can be re-derived without
re-running the LLM. Cluster wall-clock numbers reported in the paper are
the median across three independent re-runs of a 3-query subset of the
official qrels per (model, dataset, method) configuration to absorb shared-
cluster GPU-utilization variance.
